%--------------------------------------------------------------
\documentclass[ja=standard,xelatex]{bxjsarticle}
%--------------------------------------------------------------

%--------------------------------------------------------------
%packages%{
% #region
\usepackage{bigints}
\usepackage{changepage}
\usepackage[utf8]{inputenc}
\usepackage{lipsum}
\usepackage[bottom]{footmisc}
\usepackage{subcaption}
\usepackage{zxjatype}
\usepackage[hiragino]{zxjafont} 
\usepackage{tcolorbox}
\usepackage{amsmath}
\usepackage{nccmath}
\usepackage{amsthm}
\usepackage{empheq}
\usepackage{amssymb}
\usepackage{bm}
\usepackage{amsfonts}
\usepackage{graphicx}
\usepackage{tocloft}
\usepackage{ascmac}
\usepackage{latexsym}
\usepackage[version=4]{mhchem}
\usepackage{listings, xcolor}
\usepackage{mathtools}
\usepackage[colorlinks=true, linkcolor=blue , citecolor = blue]{hyperref}
\usepackage{array}
\usepackage{url}
\usepackage[english]{babel}
\usepackage[square,numbers]{natbib}
\usepackage{physics}
\usepackage{caption}
\usepackage{ulem}
\usepackage[section]{placeins}
\usepackage{xeCJKfntef}
\usepackage{longtable}
\usepackage{float}
\usepackage{multicol}
\usepackage{mathcomp}
\usepackage{IEEEtrantools}
\usepackage{flafter}
\usepackage{setspace}
% #endregion
%--------------------------------------------------------------%}

%--------------------------------------------------------------
%define%{
% #region

%%skip,space関係
\newcommand{\vs}[1]{\vskip#1\baselineskip}
\newcommand{\parlskip}[1]{\par\setlength{\leftskip}{#1pt}}

%equation環境
\newcommand{\eqn}[1]{\begin{equation*}{#1}\end{equation*}}
\newcommand{\neqn}[1]{\begin{equation}{#1}\end{equation}}
\newcommand{\leqn}[2]{\begin{fleqn}[#1pt]\begin{equation*}{#2}\end{equation*}\end{fleqn}[#1pt]}
\newcommand{\leqnarray}[2]{\begin{fleqn}[#1pt]\begin{eqnarray*}{#2}\end{eqnarray*}\end{fleqn}}
\newcommand{\eqns}[2]{\begin{equation*}\begin{array}{#1}#2\end{array}\end{equation*}}
\newcommand{\neqns}[2]{\begin{IEEEeqnarray}{#1}#2\end{IEEEeqnarray}}
\newcommand{\neqnarray}[1]{\begin{eqnarray}#1\end{eqnarray}}

%定理・定義・証明等の記述（数学）
\newcommand{\setting}[1]{\vs{0.3}\underline{\textbf{Setting}}\vs{0.3}\parlskip{20}{#1}\parlskip{0}}
\newcommand{\thm}[1]{\vs{0.3}\underline{\textbf{Thm}}\vs{0.3}\parlskip{20}{#1}\parlskip{0}}
\newcommand{\Def}[1]{\vs{0.3}\underline{\textbf{Def}}\vs{0.3}\parlskip{20}{#1}\parlskip{0}}
\newcommand{\rem}[1]{\vs{0.5}\large{\underline{{\textit{remark}}}}\normalsize\vs{0.5}{#1}}
\newcommand{\Proof}[1]{\vs{0.5}\large{\underline{\textit{Proof}}}\normalsize\vs{0.5}{#1}\setcounter{equation}{0}}
\newcommand{\free}[2]{\vs{0.5}\large{\underline{\textit{#1}}}\normalsize\vs{0.5}{#2}\setcounter{equation}{0}}

%tinysections
\newcommand{\tinysection}[3]{\vs{0.5}\large\underline{\textbf{(#1)：#2}}\normalsize\vs{0.5}{#3}}
\newcommand{\tinytinysection}[4]{\vs{0.3}\normalsize{\underline{\textbf{(#1.#2)：#3}}}\vs{0.3}{#4}}
\newcommand{\tinytinytinysection}[5]{\vs{0.2}\normalsize{\underline{(#1.#2.#3)：#4}}\vs{0.2}{#5}}

%微分
\newcommand{\dif}[2]{\dfrac{d#1}{d#2}}
\newcommand{\difop}[1]{\dfrac{d}{d#1}}
\newcommand{\lagdif}[2]{\frac{D#1}{D#2}}

%偏微分
\newcommand{\pdif}[2]{\dfrac{\partial{#1}}{\partial#2}}
\newcommand{\pdifop}[1]{\dfrac{\partial}{\partial#1}}
\newcommand{\pdifl}[1]{\partial_{#1}}
\newcommand{\pdiftwice}[2]{\dfrac{\partial^{2}{#1}}{\partial {#2} ^{2}}}

%積分
\newcommand{\dint}{\displaystyle\int}
\newcommand{\diint}{\displaystyle\iint}
\newcommand{\diiint}{\displaystyle\iiint}

%総和
\newcommand{\sumn}[2]{\sum_{n=#1}^{#2}}
\newcommand{\sumi}[2]{\sum_{i=#1}^{#2}}
\newcommand{\sumj}[2]{\sum_{j=#1}^{#2}}
\newcommand{\sumk}[2]{\sum_{k=#1}^{#2}}
\newcommand{\suml}[2]{\sum_{l=#1}^{#2}}
\newcommand{\summ}[2]{\sum_{m=#1}^{#2}}
\newcommand{\sump}[2]{\sum_{p=#1}^{#2}}
\newcommand{\sumq}[2]{\sum_{q=#1}^{#2}}
\newcommand{\sumfr}[3]{\sum_{#1=#2}^{#3}}
\newcommand{\dsum}{\displaystyle\sum}

%ベクトル
\newcommand{\vect}[1]{\mathbf{#1}}
\newcommand{\vecp}[2]{\mathbf{#1}^{#2}}
\newcommand{\vecn}[1]{\mathbf{#1}^n}
\newcommand{\vecm}[1]{\mathbf{#1}^m}
\newcommand{\pvect}[1]{\begin{pmatrix}#1\end{pmatrix}}

%括弧
\newcommand{\bracketnormal}[1]{\left(#1\right)}
\newcommand{\curlybracket}[1]{\left\{#1\right\}}

%ε-δ論法
\newcommand{\epsilondelta}[4]{\forall\;\epsilon>0,\exists\delta\;\; s.t, \;\; 0<\left| #1-#2 \right|<\delta \Rightarrow \left| #3(#1)-#4\right|<\epsilon}

%その他
\newcommand{\realpower}[1]{\mathbb{R}^{#1}}
\newcommand{\Real}{\mathbb{R}}
\newcommand{\Bf}[1]{\mathbf{#1}}
\newcommand{\limit}[2]{\lim_{#1 \to #2}}
\newcommand{\absdefined}[1]{\left| #1 \right|}
\newcommand{\map}[5]{#1:#2\to#3;\;\;#4\mapsto#5}
\newcommand{\inner}[1]{\left\langle #1 \right\rangle}
\newcommand{\vertline}[2]{\vs{0}\begin{center}\rule{#1\textwidth}{#2pt}\end{center}\vs{0}}
\newcommand{\coord}[2]{\left(#1\;,\;#2\right)}
\newcommand{\scipow}[1]{\times 10^{#1}}
\newcommand{\labeqn}[1]{\label{eqn:#1}}
\newcommand{\labfig}[1]{\label{fig:#1}}
\newcommand{\refeqn}[1]{（\ref{eqn:#1}）}
\newcommand{\reffig}[1]{Figure\ref{fig:#1}}

% #endregion
%--------------------------------------------------------------%}

%--------------------------------------------------------------
%title,author,etc%{
% #region
\title{地球物理数値解析レポート課題（天野教員担当分）}
\author{05-242628三田村彰大}
\date{}
% #endregion
%--------------------------------------------------------------%}

\begin{document}

%--------------------------------------------------------------
% #region
%code_area_settings%{
\lstset{
    basicstyle = {\ttfamily\small}, % 基本的なフォントスタイル
    frame = {tbrl}, % 枠線の枠線。t: top, b: bottom, r: right, l: left
    breaklines = true, % 長い行の改行
    numbers = left, % 行番号の表示。left, right, none
    showspaces = false, % スペースの表示
    showstringspaces = false, % 文字列中のスペースの表示
    showtabs = false, %　タブの表示
    keywordstyle = \color{blue}, % キーワードのスタイル。intやwhileなど
    commentstyle = {\color[HTML]{1AB91A}}, % コメントのスタイル
    identifierstyle = \color{black}, % 識別子のスタイル　関数名や変数名
    stringstyle = \color{brown}, % 文字列のスタイル
    captionpos = t % キャプションの位置 t: 上、b: 下
}
\captionsetup*[lstlisting]{labelformat=empty}
% #endregion
%--------------------------------------------------------------%}

%--------------------------------------------------------------
%document_settings%{
% #region
\setcounter{tocdepth}{3}
\setlength{\parindent}{0pt}
\setlength{\abovedisplayskip}{15pt}
\setlength{\belowdisplayskip}{15pt}
\setlength{\abovedisplayshortskip}{15pt}
\setlength{\belowdisplayshortskip}{15pt}
\setlength{\baselineskip}{15pt}
\newgeometry{top=25truemm,bottom=25truemm,left=20truemm,right=20truemm}
\renewcommand{\arraystretch}{1.5}
\bibliographystyle{abbrvnat}
\setlength{\columnseprule}{1pt}
%asm_settings
\theoremstyle{definition}
\newtheorem{prop}{proposition}[section]
%subfigure_settings
\captionsetup*[subfigure]{position=bottom}
\captionsetup*[figure]{labelfont=bf}
\captionsetup*[table]{labelfont=bf}
%footnote_settings
\setlength{\footnotesep}{0.4\baselineskip} % 脚注どうしの間隔
\makeatletter
\renewcommand\@makefntext[1]{%
  % ---- 脚注本文の段落設定（同一脚注内） ----
  \setlength{\parindent}{0pt}%
  \setlength{\parskip}{0.4\baselineskip}%
  % ---- ぶら下げ＋マーク表示 ----
  \noindent
  \hb@xt@1.8em{\hss\@makefnmark}%
  \hspace{0.4em}%
  \hangindent=2.2em%
  \hangafter=1%
  #1%
}
\renewcommand{\footnoterule}{%
  \kern -3pt
  \hrule width 1.0\linewidth
  \kern 5pt
}
\makeatother
\makeatletter
\renewcommand{\bibsection}{} %"References" 見出しを無効化
\makeatother
\makeatletter
\renewenvironment{proof}[1][\proofname]{%
  \par\pushQED{\qed}\normalfont
  \trivlist
  \item[\hskip\labelsep\bfseries #1\@addpunct{.}]%
  \ignorespaces
}{%
  \popQED\endtrivlist\@endpefalse
}
\makeatother
% #endregion
%--------------------------------------------------------------%}

%--------------------------------------------------------------
%maketitle_makehedding_tableofcontents%{
% #region
\maketitle

\hrule

\begin{description}[
  labelwidth=4.5cm,
  labelsep=0.5cm,
  leftmargin=5.0cm,
  style=multiline
]
  \item[数値計算結果の動画置き場：] \href{https://shota-mitamura.com/works/geo_analysiskadai1/kadai/}{https://shota-mitamura.com/works/geo_analysiskadai1/kadai/}
  \item[コード置き場：] \href{https://github.com/ramutami/modeling_mitamura}{https://github.com/ramutami/modeling\_mitamura}
  \item[AIとのチャットログ：] \href{https://chatgpt.com/g/g-p-69fabe32fd9c81919d69b22d4294ca7e-konhiyutakurahuikusulun-ke-ti/project}{https://chatgpt.com/g/g-p-69fa...}
\end{description}

\hrule
\vs{1}


\tableofcontents

% #endregion
%--------------------------------------------------------------%}

%start_document
\section{課題A}
\section{課題B}
\begin{spacing}{2.0}
（B-1）$u\bracketnormal{x,t} = u_0\bracketnormal{x-tc}$　（B-2）$\pdif{u}{t} = \dfrac{u_i^{n+1}-u_i^n}{\Delta t}$　（B-3）$\pdif{u}{x}=\dfrac{u_{i+1}^n-u_{i-1}^n}{2\Delta x}$　
\vs{0}
（B-4）$u_i^{n+1} = u_i^n - \dfrac{c\Delta t}{2\delta x}\bracketnormal{u_i^{n+1}-u_{i-1}^n} $　（B-5）無条件不安定　
\vs{0}（B-6）$0\leq \nu \leq 1$なら安定、ただし$\nu = \dfrac{c\Delta t}{\Delta x}$　（B-7）$-1\leq \nu \leq 0$なら安定、ただし$\nu = \dfrac{c\Delta t}{\Delta x}$　
\vs{0}
\end{spacing}
（B-8）数値計算スキームが安定となるためには、物理的な情報の伝播がすでに済んでいる位置での情報を使って求めたい位置での時間ステップ更新を行う必要がある。
\begin{spacing}{2.0}
（B-9）位相誤差　（B-10）位相速度のズレが大きくなる　（B-11）$\dfrac{1}{2}c^2 \pdiftwice{u}{x}\Delta t$　（B-12）逆拡散　
\vs{0}
（B-13）数値拡散の項（あるいは、$\curlybracket{\dfrac{1}{2}\nu ^2\bracketnormal{u_{i+1}^n - 2u_{i}^n+u_{i-1}^n}}$の項）　（B-14）数値振動　
\vs{0}
（B-15）Godunovの定理　（B-16）人工粘性　（B-17）積分　（B-18）保存則が成立　（B-19）特性量　
\vs{0}
（B-20）$\dif{x}{t}=\lambda\bracketnormal{x}$　（B-21）線形移流方程式が$\pdif{u}{t}+c\pdif{u}{x}$と書き表される時、移流速度$c$　（B-22）物理量$u$　
\vs{0}
（B-23）２つの特性量の値を持つことが必要　（B-24）衝撃波　（B-25）粘性（数値拡散）　（B-26）弱解　
\vs{0}
（B-27）エントロピー条件　（B-28）左向き（負向き）の音波　（B-29）エントロピー波　（B-30）右向き（正向き）の音波
\end{spacing}

\section{課題D-1}
　$x>0$向きの移流速度をもつ線形移流方程式を考えることにする。すなわち次の方程式を考える。
\neqn{
  \pdif{u}{t} + c\pdif{u}{x} = 0 \quad \bracketnormal{c>0}
}
ただし$c$は計算に際して適宜定めるものとする。
\vs{0}
　また、今回は一次元風上差分とLax-Wendroffを用いて上式を解くことを考える。具体的には、空間方向に幅$\Delta x$で$i=1,2,\ldots$と離散化し、時間方向に幅$\Delta t$で$n=0,1,\ldots$と離散化する時、一次元風上差分及びLax-Wendroffは（$c>0$よりFTBSとなるので）次のようなスキームとなる。
\neqns{ll}{
  \text{FTBS:}\quad & u_i^{n+1} = u_i^n - \dfrac{c\Delta t}{\Delta x}\bracketnormal{u_i^n - u_{i-1}^{n}}\\[5pt]
  \text{Lax-Wendroff:}\quad & u_i^{n+1} = u_i^n - \dfrac{c\Delta t}{2\Delta x}\bracketnormal{u_{i+1}^n - u_{i-1}^{n}} + \dfrac{1}{2}\bracketnormal{\dfrac{c\Delta t}{\Delta x}}^2 \bracketnormal{u_{i+1}^n - 2u_i^n + u_{i-1}^n}
}
　実際の計算においては、$c=0.5$と定め、適当に定めた格子幅$\Delta x$に対しCFL条件を満たすように（すなわち$c\Delta t/\Delta x <1 $となるように）$\Delta t$を定めた。\footnote{
  具体的には、$\Delta x=0.01$、$c\Delta t/\Delta x = 1/3$とした。
}


\section{課題D-2}

\end{document}


